\documentclass[10pt,a4paper]{amsart}
\usepackage[margin=19mm]{geometry}
\usepackage{amsmath,amssymb,mathtools}
\usepackage[scr=rsfs,cal=euler]{mathalfa}
\usepackage{xfrac}
\usepackage[unicode,hidelinks,bookmarks=false]{hyperref}
\newcommand{\ie}{i.e.\ }
\newcommand{\cf}{cf.\ }
\newcommand{\resp}{respectively\ }
\newcommand{\Subsection}{\subsection}
\providecommand{\nobreakdash}{\nobreak\hbox{-}}
\setlength{\emergencystretch}{2em}
\multlinegap0pt
\usepackage{amssymb}
\usepackage{mathtools}
\newtheorem{lemma}{Lemma}
\newtheorem{theorem}{Theorem}
\newcommand{\BV}{\text{BV}}
\newcommand{\R}{\mathbb{R}}
\title{Anisotropic gradient rearrangement of BV functions and applications}
\author{Gloria Paoli, Yabo Yang}
\date{Source: arXiv:2605.15849v1 (CC BY 4.0)}
\begin{document}
\maketitle

\noindent\emph{Source and licence.} Anisotropic gradient rearrangement of BV functions and applications. Version arXiv:2605.15849v1 (CC BY 4.0), distributed by arXiv under the Creative Commons Attribution 4.0 International licence, http://creativecommons.org/licenses/by/4.0/, as recorded in the arXiv licence metadata for this exact version and shown on the abstract page https://arxiv.org/abs/2605.15849v1. Full licence notice: \texttt{licenses/arxiv-2605.15849-CC-BY-4.0.txt}.\par\smallskip

\noindent\textit{Editorial scope.} Counted unit: the article's theorem thm:main\_comparison (source0.tex lines 653--661). The article's complete original proof of it is delivered with it (source0.tex lines 664--887, one \texttt{proof} environment of 224 lines). No mathematics is written, repaired or re-derived: the statement and the proof are the article's own lines, in the article's own notation, and no retained line is substituted or re-flowed.  Retained uncounted as the source-local context the proof needs: lines 337--339; lines 424--429; lines 446--470; lines 506--509; lines 511--517; lines 541--547; lines 556--564.  Nothing was omitted from any retained span, so the package carries no omission record.  No retained line contains a citation, so the wrapper carries no bibliography and no bibliography entry.  No retained line contains a figure environment, an image-inclusion command or drawing code, so no image is delivered and the package declares no asset dependency.  Editorial formatting only, in addition: an AMS \texttt{amsart} preamble that restates the article's own declarations and macros used by the retained lines, namely the environments \texttt{lemma}, \texttt{theorem} on separate counters, together with the standard packages those lines need.  The retained environments print in this document as Lemma 1, Theorem 1 in order of appearance, whereas the article prints the same items with the numbering of its own full document.  The complete native source of the article as distributed in the arXiv e-print for this exact version is bundled unchanged as \texttt{sources/200686/source0.tex}.
\begin{equation} \label{eq:isoperimetric}
    P_H(E) \geq n \kappa_n^{1/n} |E|^{1 - 1/n}.
\end{equation}

\begin{lemma}\label{lem:G_derivative}
Let $u$ be a non-negative function in $\BV_0(\Omega)$. Then, for a.e. $s\in[0,+\infty)$,
\begin{equation}\label{eq:G_integral}
    G(s)=\int_{u^*(s)}^{+\infty} P_H(\{u>t\})\,dt.
\end{equation}
\end{lemma}

\begin{lemma}\label{lem:approximation}
Let \(u\in \BV_0(\Omega)\) be non-negative and let \(g\in L^1(\Omega)\). Define
\[
G_g(s):=\int_{\{u>u^*(s)\}} g(x)\,dx,
\qquad s\in[0,|\Omega|].
\]
Assume that \(G_g\) is absolutely continuous and write
\[
G_g(s)=\int_0^s F_g(\tau)\,d\tau .
\]
Then there exists a sequence \(\{g_k\}_{k\in\mathbb N}\subset L^1(0,|\Omega|)\) such that
\[
g_k^*=g^*
\qquad\text{for every }k,
\]
and
\[
\lim_{k\to\infty}
\int_0^{|\Omega|}g_k(\tau)\varphi(\tau)\,d\tau
=
\int_0^{|\Omega|}F_g(\tau)\varphi(\tau)\,d\tau
\qquad
\forall\,\varphi\in \BV([0,|\Omega|)).
\]
\end{lemma}

\begin{equation}\label{eq:sigma_decomposition}
d\sigma_H = F_{H,1}(s)\,ds + d\sigma_{H,2}(s)
\qquad \text{on }[0,|\Omega|].
\end{equation}

\begin{equation}\label{eq:sigma2_total_mass}
\int_{(0,|\Omega|]} d\sigma_{H,2}(\tau)
=
G_{H,2}(|\Omega|)
=
|D^s u|_H(\mathbb{R}^n).
\end{equation}

\begin{equation}\label{eq:ustar_problem}
\begin{cases}
H(\nabla u^\star)(x)=g_\sharp(x) & \text{for a.e. }x\in\Omega^\star,\\[1ex]
u^\star \text{ is } H^\circ\text{-radial and nonincreasing},\\[1ex]
u^\star(x)=c_0 & \text{for }x\in\partial\Omega^\star,
\end{cases}
\end{equation}

\begin{equation}\label{eq:ustar_rearrangement_explicit}
(u^\star)^*(s)
=
\int_s^{|\Omega|}
\frac{g_*(t)}{n\kappa_n^{1/n} t^{1-\frac1n}}\,dt
+
\frac{|D^s u|_H(\mathbb{R}^n)}{n\kappa_n^{1/n}|\Omega|^{1-\frac1n}},
\qquad s\in[0,|\Omega|].
\end{equation}

\begin{theorem}\label{thm:main_comparison}
Let $\Omega\subset\R^n$ be a bounded open set with finite perimeter, and let
$\Omega^\star$ be the centered Wulff ball such that $|\Omega^\star|=|\Omega|$.
Assume that $u$ is a non-negative function in $\BV_0(\Omega)$. If $u^\star$ is the
anisotropic gradient rearrangement defined by \eqref{eq:ustar_problem}, then
\begin{equation}\label{eq:main_inequality}
\|u\|_{L^1(\Omega)}\le \|u^\star\|_{L^1(\Omega^\star)}.
\end{equation}
\end{theorem}

\begin{proof}
Let $u^*$ be the decreasing rearrangement of $u$. By Lemma~\ref{lem:G_derivative}, for a.e.
$s\in[0,|\Omega|]$,
\[
G(s)=\int_{u^*(s)}^{+\infty} P_H(\{u>t\})\,dt .
\]
 Let
$0\le s_1<s_2\le |\Omega|$. Since \(G\) is increasing and
\[
G(s)=\sigma_H((0,s]),
\]
we have
\[
\sigma_H((s_1,s_2))
=
G(s_2^-)-G(s_1),
\]
where \(G(s_2^-):=\lim_{s\to s_2^-}G(s)\). Using the representation of \(G\)
and passing to the left limit, we obtain
\begin{align}
\sigma_H((s_1,s_2))
&=
\lim_{s\to s_2^-}
\left[
G(s)-G(s_1)
\right] \notag\\
&=
\lim_{s\to s_2^-}
\int_{u^*(s)}^{u^*(s_1)} P_H(\{u>t\})\,dt . \label{eq:sigma_interval_left}
\end{align}
By the anisotropic isoperimetric inequality \eqref{eq:isoperimetric}
\begin{align}
\sigma_H((s_1,s_2))
&\ge
\lim_{s\to s_2^-}
\int_{u^*(s)}^{u^*(s_1)}
n\kappa_n^{1/n}\mu(t)^{1-\frac1n}\,dt .
\label{eq:interval_measure_lower}
\end{align}
Define now 
\[
\mathcal K(\tau):=
\int_\tau^{+\infty}
n\kappa_n^{1/n}\mu(t)^{1-\frac1n}\,dt,
\qquad \tau\ge 0.
\]
The function \(\mathcal K\) is finite and Lipschitz continuous. Indeed, by the
anisotropic isoperimetric inequality and the anisotropic coarea formula,
\[
\int_0^{+\infty}
n\kappa_n^{1/n}\mu(t)^{1-\frac1n}\,dt
\le
\int_0^{+\infty} P_H(\{u>t\})\,dt
=
|Du|_H(\mathbb R^n)<+\infty,
\]
and the integrand is bounded by
\(n\kappa_n^{1/n}|\Omega|^{1-\frac1n}\).
With this notation, \eqref{eq:interval_measure_lower} gives
\[
\sigma_H((s_1,s_2))
\ge
D[\mathcal K(u^*)]((s_1,s_2)).
\]
By the regularity of Radon measures, this implies
\begin{equation}\label{eq:measure_domination}
\sigma_H(E)\ge D[\mathcal K(u^*)](E)
\qquad
\text{for every Borel set }E\subset(0,|\Omega|).
\end{equation}

We now consider \(D[\mathcal K(u^*)]\). Since \(u^*\) is monotone decreasing,
\(-Du^*\) is a positive measure. By the chain rule for one-dimensional BV
functions,
\begin{equation}\label{eq:chain_rule_K}
D[\mathcal K(u^*)]
=
n\kappa_n^{1/n}s^{1-\frac1n}\,(-Du^*)
\qquad\text{ on }(0,|\Omega|).
\end{equation}
 Indeed, on the
absolutely continuous and Cantor parts, one uses
\[
\mu(u^*(s))=s
\qquad\text{for }|Du^*|\text{-a.e. }s.
\]
At a jump point \(s\in J_{u^*}\), one has
\[
\mu(t)=s
\qquad
\text{for }t\in\big((u^*)_+(s),(u^*)_-(s)\big).
\]
Hence
\begin{align*}
D[\mathcal K(u^*)](\{s\})
&=
\mathcal K((u^*)_+(s))-\mathcal K((u^*)_-(s))\\
&=
\int_{(u^*)_+(s)}^{(u^*)_-(s)}
n\kappa_n^{1/n}\mu(t)^{1-\frac1n}\,dt\\
&=
n\kappa_n^{1/n}s^{1-\frac1n}
\big((u^*)_-(s)-(u^*)_+(s)\big)\\
&=
n\kappa_n^{1/n}s^{1-\frac1n}(-Du^*)(\{s\}).
\end{align*}
This proves \eqref{eq:chain_rule_K}.

Combining \eqref{eq:measure_domination} and \eqref{eq:chain_rule_K}, we get
\[
n\kappa_n^{1/n}s^{1-\frac1n}\,(-Du^*)
\le
d\sigma_H
\qquad\text{on }(0,|\Omega|).
\]
Since \(u^*(|\Omega|)=0\), it follows that, for a.e. \(s\in(0,|\Omega|)\),
\begin{equation}\label{eq:u_pointwise_estimate}
u^*(s)
\le
\int_{(s,|\Omega|]}
\frac{d\sigma_H(\tau)}
{n\kappa_n^{1/n}\tau^{1-\frac1n}}.
\end{equation}
Integrating \eqref{eq:u_pointwise_estimate} over $[0,|\Omega|]$ and using Fubini's
theorem, we get
\begin{align}
\|u\|_{L^1(\Omega)}
&=
\int_0^{|\Omega|} u^*(s)\,ds \notag\\
&\le
\int_0^{|\Omega|}
\left(
\int_s^{|\Omega|}
\frac{d\sigma_H(\tau)}{n\kappa_n^{1/n}\tau^{1-\frac1n}}
\right)ds \notag\\
&=
\frac1{n\kappa_n^{1/n}}
\int_0^{|\Omega|}\tau^{\frac1n}\,d\sigma_H(\tau). \label{eq:L1_after_fubini}
\end{align}
Using the decomposition \eqref{eq:sigma_decomposition}, this becomes
\begin{equation}\label{eq:L1_split}
\|u\|_{L^1(\Omega)}
\le
\frac1{n\kappa_n^{1/n}}
\left[
\int_0^{|\Omega|}\tau^{\frac1n}F_{H,1}(\tau)\,d\tau
+
\int_{(0,|\Omega|]}\tau^{\frac1n}\,d\sigma_{H,2}(\tau)
\right].
\end{equation}

We now estimate the first term. Set
\[
g:=H(\nabla^a u).
\]
By  Lemma \ref{lem:approximation}, applied to \(g\) with
respect to the level sets of \(u\), there exists a sequence \(\{g_k\}\subset L^1(0,|\Omega|)\)
such that
\[
g_k^*=g^*
\]
and
\[
\lim_{k\to\infty}
\int_0^{|\Omega|} g_k(\tau)\varphi(\tau)\,d\tau
=
\int_0^{|\Omega|}F_{H,1}(\tau)\varphi(\tau)\,d\tau
\qquad
\forall\varphi\in BV([0,|\Omega|)).
\]
Choosing \(\varphi(\tau)=\tau^{1/n}\) and using the Hardy--Littlewood inequality,
we obtain
\[
\int_0^{|\Omega|}\tau^{1/n}F_{H,1}(\tau)\,d\tau
\le
\int_0^{|\Omega|}\tau^{1/n}g_*(\tau)\,d\tau
=
\int_0^{|\Omega|}\tau^{1/n}\big(H(\nabla^a u)\big)_*(\tau)\,d\tau .
\]
For the second term, since $\tau^{1/n}\le |\Omega|^{1/n}$ on $[0,|\Omega|]$, by
\eqref{eq:sigma2_total_mass} we have
\begin{equation}\label{eq:sigma2_estimate}
\int_{(0,|\Omega|]}\tau^{\frac1n}\,d\sigma_{H,2}(\tau)
\le
|\Omega|^{\frac1n}|D^s u|_H(\mathbb{R}^n).
\end{equation}
Combining this with \eqref{eq:L1_split}, and \eqref{eq:sigma2_estimate},
we obtain
\begin{equation}\label{eq:L1_upper_bound}
\|u\|_{L^1(\Omega)}
\le
\frac1{n\kappa_n^{1/n}}
\int_0^{|\Omega|}\tau^{\frac1n} g_*(\tau)\,d\tau
+
\frac{|\Omega|^{1/n}}{n\kappa_n^{1/n}}|D^s u|_H(\mathbb{R}^n).
\end{equation}

It remains to compute the $L^1$ norm of $u^\star$. By
\eqref{eq:ustar_rearrangement_explicit},
\[
(u^\star)^*(s)
=
\int_s^{|\Omega|}
\frac{g_*(t)}{n\kappa_n^{1/n}t^{1-\frac1n}}\,dt
+
\frac{|D^s u|_H(\mathbb{R}^n)}{n\kappa_n^{1/n}|\Omega|^{1-\frac1n}}.
\]
Integrating over $[0,|\Omega|]$ and using Fubini's theorem once again, we get
\begin{align*}
\|u^\star\|_{L^1(\Omega^\star)}
&=
\int_0^{|\Omega|}(u^\star)^*(s)\,ds \\
&=
\frac1{n\kappa_n^{1/n}}
\int_0^{|\Omega|} t^{\frac1n} g_*(t)\,dt
+
\frac{|\Omega|^{1/n}}{n\kappa_n^{1/n}}|D^s u|_H(\mathbb{R}^n).
\end{align*}
Comparing this identity with \eqref{eq:L1_upper_bound}, we conclude that
\[
\|u\|_{L^1(\Omega)}\le \|u^\star\|_{L^1(\Omega^\star)}.
\]
This proves \eqref{eq:main_inequality}.
\end{proof}

\end{document}
